% ---
% id: statistics/binomial-box/lem:cheb
% kind: lem
% title: "Chebyshev's inequality for a binomial count"
% summary: >-
%   A binomial (n, p) count deviates from np by at least t > 0 with
%   probability at most np(1 − p)/t². This gives the finite-sample width bound
%   of the exact rule in [[thm:monotonicity-limits]] (iv).
% prominence: technical
% status: accepted
% uses: []
% assumes: []
% refs: []
% provenance: newly-derived
% lean:
% ---

\begin{lemma}\label{lem:cheb}
Let $\Ncount{}$ be binomial with parameters $(n,p)$ and let $t>0$. Then
$\Pr\big[|\Ncount{}-np|\ge t\big]\le np(1-p)/t^{2}$.
\end{lemma}

\begin{proof}
Write $\Ncount{}=\sum_{j=1}^{n}\chi_j$ with $\chi_j$ independent
Bernoulli$(p)$; then $\mathbb{E}[\Ncount{}]=np$ and, by independence,
$\operatorname{Var}(\Ncount{})=\sum_{j}\operatorname{Var}(\chi_j)=np(1-p)$.
Consider the non-negative random variable $(\Ncount{}-np)^{2}$. Pointwise it is
at least $t^{2}$ wherever $|\Ncount{}-np|\ge t$ and at least $0$ elsewhere, so
it dominates the random variable taking the value $t^{2}$ wherever
$|\Ncount{}-np|\ge t$ and $0$ elsewhere. Taking expectations,
$t^{2}\,\Pr\big[|\Ncount{}-np|\ge t\big]
\le\mathbb{E}\big[(\Ncount{}-np)^{2}\big]
=\operatorname{Var}(\Ncount{})=np(1-p)$, which is the claim.
\end{proof}
